\subsection{A compact four-node degeneration}
\label{subsec:x9-mirror-periods}

Use the primitive nef divisors $H_a$ and algebraic mirror
coordinates $z_a$ specified in Appendix~\ref{app:mirror-periods}.
Their dual integral curve classes are $\beta_a$, with
\begin{equation}
 C_2=\beta_1,\qquad C_1=\beta_5,\qquad
 \iota_*e_2=\beta_6,\qquad
 \iota_*e_1=\beta_1+\beta_5+\beta_6.
 \label{eq:x9-electric-nef-charges}
\end{equation}
The notation $\beta_a$ for curve classes is distinct from the
local moment-map coordinate $\beta$. Set
\begin{equation}
 \Sigma:\ z_1=z_5=z_6=1,\qquad (v,a,b)=(z_2,z_3,z_4).
 \label{eq:x9-finite-bulk-locus}
\end{equation}
In a smooth ambient chart $(x,y,r,s)\in(\CC^*)^2\times\CC^2$
the full compact mirror equation restricts to
\begin{equation}
 \mathcal F=(1+xy)(1+x^{-1}+vy^{-1}+r)
             +s\bigl(r(br-1)+as\bigr)=0.
 \label{eq:x9-finite-bulk-factorization}
\end{equation}
This is not a truncation in $a$ or $b$. Its four nodes lie at
\begin{equation}
 s=0,\quad y=-x^{-1},\quad r\in\{0,b^{-1}\},\quad
 vx^2-(1+r)x-1=0.
 \label{eq:x9-finite-bulk-nodes}
\end{equation}
For $bv\ne0$ and $(1+r)^2+4v\ne0$, the Hessian determinants
are nonzero. The two nodes at $r=0$ approach the local $E_2$
doublet; the other pair approaches the original conifold sectors
in their square charts as $b\to0$. Appendix~\ref{app:mirror-periods}
constructs the charts and excludes additional singularities at
an exact rational sample, giving an open four-node stratum.

The local smoothing covectors in the three normal parameters
$(z_1,z_5,z_6)$ are
\begin{equation}
 (x^{-1},x^{-1}+r,-1).
 \label{eq:x9-finite-bulk-smoothing}
\end{equation}
They have rank three. The closure of $s=1+xy=0$ is a compact
Weil divisor passing once through each node. Its complement of
the node neighborhoods supplies a four-chain with boundary the
four oriented vanishing spheres. Thus there is one relation with
unit coefficients, and the smoothing rank excludes a second.
Blowing up this divisor gives a projective small resolution.

The four vanishing cycles are mutually local, with rank-three
span. Standard type-IIB conifold physics therefore supplies four
massless hypermultiplets charged under three Abelian vectors
\cite{Strominger:1995Conifold,Greene:1995hu}. Their Higgsing gives
$(\Delta n_V,\Delta n_H^{\rm neutral})=(-3,+1)$ and mirror
Hodge numbers $(122,6)\to(123,3)$. This is the finite-radius
particle description. The family identification below establishes
more than this numerical agreement with the original smoothing.

\subsection{The integral massless charges}

\begin{proposition}\label{prop:x9-integral-charges}
On a marked sufficiently small nonzero $(v,a,b)$ branch of $\Sigma$,
the large-volume flat coordinates satisfy
\begin{equation}
 t_1=t_5=t_6=0.
 \label{eq:x9-finite-bulk-flat-locus}
\end{equation}
In D6--D4--D2--D0 order, the four vanishing charges are
\begin{equation}
 (0,0,C_1,0),\quad(0,0,C_2,0),\quad
 (0,0,\iota_*e_1,0),\quad(0,0,\iota_*e_2,0).
 \label{eq:x9-finite-bulk-integral-charges}
\end{equation}
Their span $V$ is a primitive rank-three isotropic sublattice.
Its unique relation, in this order, is $(-1,-1,1,-1)$.
\end{proposition}
\begin{proof}
At fixed spectator degree, summation over all three local degrees
gives convergent expressions for the scalar period $W_0$ and
the six electric periods $W_0t_a$. A binomial convolution sets
the three local logarithmic corrections to zero on $\Sigma$.
The branch is continued without coordinate winding, fixing the
integer logarithmic offsets. The estimates and coefficient
identities are proved in Appendix~\ref{app:mirror-periods}.

The seven electric periods are single-valued near the locus.
Picard--Lefschetz monodromy around each individual node therefore
excludes D6 and D4 components from its vanishing charge. The
three spectator flat coordinates are independent along $\Sigma$,
excluding spectator D2 components and a D0 shift. Unit-normalized
node residues, matched to the exceptional-curve and nodal flop
limits, fix the four primitive labels in
\eqref{eq:x9-electric-nef-charges}. The first, fifth and sixth
dual basis elements generate their span, proving saturation;
the integral cohomology tests in the appendix exclude hidden
torsion. Mutual locality follows either from their pure D2
charges or from their disjoint Milnor neighborhoods.
\end{proof}

The massless charge matrix, after removing neutral vectors, is
\begin{equation}
 Q_{\rm light}=\begin{pmatrix}1&0&1&0\\0&1&1&0\\0&0&1&1\end{pmatrix}.
 \label{eq:x9-local-four-charge-matrix}
\end{equation}
Its three Smith invariants are one. The formal zero-level
hyperk\"ahler quotient is the same $\CC^2/\ZZ_4$ surface as
\eqref{eq:x9-dressed-generators}. Here all four columns have an actual
simultaneous massless-particle interpretation at finite radius.
This does not determine a compact quaternionic-K\"ahler metric.

At a generic Higgs vacuum, a magnetic divisor charge $P$ has
oriented flux vector
\[
 (-PC_1,-PC_2,P\iota_*e_1,-P\iota_*e_2).
\]
Its image is all of $\{w\in\ZZ^4:\sum_iw_i=0\}$, since
the three independent electric pairings are surjective onto
$\ZZ^3$. This is an $A_3$ winding lattice, not an $SU(4)$
enhancement. Its kernel is exactly
\eqref{eq:x9-unconfined-magnetic-lattice}; these are the magnetic
directions whose periods can descend through the transition.
The saturation also implies a connected unbroken gauge torus.

\subsection{Hypermultiplet thresholds and their monodromy}
\label{subsec:x9-conifold-thresholds}

The four charges determine the singular gauge thresholds without
using the regular-period calculation. Near a generic point of $\Sigma$,
choose transverse special coordinates $x=(t_5,t_1,t_6)$ and let
$q_\alpha$, $\alpha=1,\ldots,4$, be the columns of
\eqref{eq:x9-local-four-charge-matrix}. In the normalization $X^0=1$
their holomorphic central charges are
\begin{equation}
 (Z_1,Z_2,Z_3,Z_4)
   =(x_1,x_2,x_1+x_2+x_3,x_3),\qquad
 Z_\alpha=q_\alpha^{\mathsf T}x.
 \label{eq:x9-threshold-central-charges}
\end{equation}
At fixed finite auxiliary radius, work off the discriminant and at
energies below these four masses, the Kaluza--Klein scale and the
other massive thresholds. Integrating out the four charged
hypermultiplets gives the universal one-loop singularity of the
holomorphic prepotential \cite[Sections~2 and~4]{Strominger:1995Conifold}:
\begin{equation}
 F=F_{\rm reg}+F_{\rm sing},\qquad
 F_{\rm sing}=\frac1{4\pi i}\sum_{\alpha=1}^4
 Z_\alpha^2\left(\log\frac{Z_\alpha}{z_*}-\frac32\right).
 \label{eq:x9-threshold-prepotential}
\end{equation}
Here $z_*\ne0$ is a dimensionless reference scale, the logarithms
are continued on the chosen branch, and $F_{\rm reg}$ is holomorphic
across $\Sigma$. Changing $z_*$ or the displayed quadratic subtraction
only redistributes holomorphic quadratic terms between the two pieces;
it does not change their sum in the fixed integral frame.

Twice differentiating gives the singular holomorphic coupling matrix
\begin{equation}
 \mathcal T^{\rm sing}_{ij}
 =\frac{\partial^2F_{\rm sing}}{\partial x_i\partial x_j}
 =\frac1{2\pi i}\sum_{\alpha=1}^4
 q_{\alpha i}q_{\alpha j}\log\frac{Z_\alpha}{z_*}.
 \label{eq:x9-threshold-matrix}
\end{equation}
A positive loop around the $\alpha$th discriminant therefore gives
\begin{equation}
 \Delta_\alpha(\partial_xF)=q_\alpha Z_\alpha,\qquad
 \Delta_\alpha\mathcal T=q_\alpha q_\alpha^{\mathsf T}.
 \label{eq:x9-threshold-monodromy}
\end{equation}
These are precisely the unit Picard--Lefschetz shifts of the marked
electric charges. If $x=\epsilon\xi$ with every
$q_\alpha^{\mathsf T}\xi\ne0$, winding $\epsilon$ once winds all
four central charges. The combined shift is
\begin{equation}
 \mathsf B=Q_{\rm light}Q_{\rm light}^{\mathsf T}
 =\begin{pmatrix}2&1&1\\1&2&1\\1&1&2\end{pmatrix},\qquad
 \mathcal T^{\rm sing}
 =\frac{\log\epsilon}{2\pi i}\mathsf B+O(1).
 \label{eq:x9-threshold-total}
\end{equation}
Thus the logarithmic coefficient has rank three, although four
hypermultiplets contribute. The Smith invariants of $\mathsf B$ are
$(1,1,4)$, whereas those of $Q_{\rm light}$ are $(1,1,1)$.
The index-four image of the combined monodromy is not the primitive
vanishing lattice $V$ and does not imply a residual discrete gauge
group. Conifold reduction uses $V^\perp/V$.

The light charges have no components along the surviving coordinates
$(t_2,t_3,t_4)$. Hence the singular terms have zero tangential
derivatives on $\Sigma$: there is no logarithmic threshold for the
unbroken neutral vectors. This does not determine their regular
couplings. The remaining subsections identify those couplings with
the smooth phase by continuing the full reduced period vector.
At $Z_\alpha=0$ the integrated-out action itself is inapplicable;
the light hypermultiplets must be restored before Higgsing. Moreover,
$\mathcal T$ is a holomorphic Hessian, not the full electromagnetic
coupling matrix of compact supergravity, reconstructed below.
